% Options for packages loaded elsewhere
\PassOptionsToPackage{unicode}{hyperref}
\PassOptionsToPackage{hyphens}{url}
\PassOptionsToPackage{dvipsnames,svgnames,x11names}{xcolor}
%
\documentclass[
  11pt,
]{article}
\usepackage{amsmath,amssymb}
\usepackage{iftex}
\ifPDFTeX
  \usepackage[T1]{fontenc}
  \usepackage[utf8]{inputenc}
  \usepackage{textcomp} % provide euro and other symbols
\else % if luatex or xetex
  \usepackage{unicode-math} % this also loads fontspec
  \defaultfontfeatures{Scale=MatchLowercase}
  \defaultfontfeatures[\rmfamily]{Ligatures=TeX,Scale=1}
\fi
\usepackage{lmodern}
\ifPDFTeX\else
  % xetex/luatex font selection
\fi
% Use upquote if available, for straight quotes in verbatim environments
\IfFileExists{upquote.sty}{\usepackage{upquote}}{}
\IfFileExists{microtype.sty}{% use microtype if available
  \usepackage[]{microtype}
  \UseMicrotypeSet[protrusion]{basicmath} % disable protrusion for tt fonts
}{}
\makeatletter
\@ifundefined{KOMAClassName}{% if non-KOMA class
  \IfFileExists{parskip.sty}{%
    \usepackage{parskip}
  }{% else
    \setlength{\parindent}{0pt}
    \setlength{\parskip}{6pt plus 2pt minus 1pt}}
}{% if KOMA class
  \KOMAoptions{parskip=half}}
\makeatother
\usepackage{xcolor}
\usepackage[margin=1in]{geometry}
\usepackage{longtable,booktabs,array}
\usepackage{calc} % for calculating minipage widths
% Correct order of tables after \paragraph or \subparagraph
\usepackage{etoolbox}
\makeatletter
\patchcmd\longtable{\par}{\if@noskipsec\mbox{}\fi\par}{}{}
\makeatother
% Allow footnotes in longtable head/foot
\IfFileExists{footnotehyper.sty}{\usepackage{footnotehyper}}{\usepackage{footnote}}
\makesavenoteenv{longtable}
\setlength{\emergencystretch}{3em} % prevent overfull lines
\providecommand{\tightlist}{%
  \setlength{\itemsep}{0pt}\setlength{\parskip}{0pt}}
\setcounter{secnumdepth}{-\maxdimen} % remove section numbering
\ifLuaTeX
  \usepackage{selnolig}  % disable illegal ligatures
\fi
\IfFileExists{bookmark.sty}{\usepackage{bookmark}}{\usepackage{hyperref}}
\IfFileExists{xurl.sty}{\usepackage{xurl}}{} % add URL line breaks if available
\urlstyle{same}
\hypersetup{
  pdftitle={Conversational Domain Adaptation of IndicTrans2 across 21 Indic Languages via Experience Replay and Model Soups},
  pdfauthor={Aditya Pratap Singh (Independent Researcher)},
  colorlinks=true,
  linkcolor={blue},
  filecolor={Maroon},
  citecolor={Blue},
  urlcolor={blue},
  pdfcreator={LaTeX via pandoc}}

\title{Conversational Domain Adaptation of IndicTrans2 across 21 Indic
Languages via Experience Replay and Model Soups}
\author{Aditya Pratap Singh (Independent Researcher)}
\date{June 2026}

\begin{document}
\maketitle

\hypertarget{abstract}{%
\subsection{Abstract}\label{abstract}}

State-of-the-art multilingual machine-translation (MT) systems such as
IndicTrans2 are trained on broad, general-domain corpora; they translate
formal text well but are weaker on casual \emph{conversational}
register. We study domain adaptation of IndicTrans2-1B to conversational
English$\to$Indic translation across \textbf{21 Indic languages}, using only
publicly available conversational data (OpenSubtitles, BPCC-H-Daily,
Tatoeba). We confirm that naive fine-tuning improves conversational chrF
but induces \textbf{catastrophic forgetting} of general-domain quality
($-$3.9 chrF on FLORES for Hindi). Applying two established techniques ---
\textbf{experience replay} (mixing general-domain data back into
fine-tuning) and \textbf{model souping} (weight-averaging the fine-tuned
model with the base) --- we obtain a single model that \textbf{beats the
IndicTrans2 base on conversational chrF in all 21 languages} (mean +6.2,
range +1.3 to +12.2) while \emph{matching} it on general-domain FLORES
quality (mean change $-$0.17, all within $\pm$0.7 chrF). Paired bootstrap
resampling confirms the conversational chrF gains are significant (p $\le$
0.004) and that FLORES is not significantly degraded. We additionally
report a negative result: the lowest-resource languages remain bounded
by the base model's absolute quality, which conversational adaptation
cannot raise. All results are automatic-metric (chrF2); we are explicit
(§6, §8) about where this metric can overstate domain-adaptation gains.
We claim no methodological novelty; our contribution is the systematic,
honestly-evaluated empirical study in the Indic conversational setting.

\hypertarget{introduction}{%
\subsection{1. Introduction}\label{introduction}}

Neural MT systems are typically optimised for, and evaluated on,
general-domain benchmarks such as FLORES-200 {[}Goyal et al., 2022; NLLB
Team, 2022{]}. Yet a large fraction of real-world translation is
\emph{conversational}: chat messages, subtitles, voice-assistant
commands. This register differs systematically from news or encyclopedic
text --- sentences are shorter, more colloquial and idiomatic, and
elliptical. A model that is excellent on FLORES can still produce
stilted, overly formal output on a casual sentence.

For Indian languages, IndicTrans2 {[}Gala et al., 2023{]} is the
strongest open English$\leftrightarrow$Indic system, trained on the 230M-pair Bharat
Parallel Corpus Collection (BPCC). It is, however, a general-domain
model. We ask a practical question:

\begin{quote}
\emph{Can we adapt a strong multilingual MT model to conversational
register across many languages, without degrading its general-domain
quality, using only public data?}
\end{quote}

Naive fine-tuning on conversational data is the obvious first attempt,
but it is well known to cause \textbf{catastrophic forgetting}
{[}McCloskey and Cohen, 1989; French, 1999{]}: the model overfits the
new domain and loses competence on the old one. We show this happens
here, then mitigate it with two established techniques applied
systematically across the Indic family.

\textbf{Contributions.} 1. A systematic study of conversational
adaptation for \textbf{English$\to$Indic MT across 21 languages}, built
entirely from public conversational corpora. 2. Evidence that
\textbf{experience replay + model souping} beats the base on
conversational chrF --- across the whole family, not a single language,
with statistical-significance testing --- while preserving
general-domain chrF. 3. An \textbf{honest cost/limitation analysis}:
catastrophic forgetting under naive fine-tuning; inflated gains on easy
in-distribution test sets; the metric's own blind spots (§6); and a
negative result on the lowest-resource languages.

We emphasise that the techniques we use --- model soups {[}Wortsman et
al., 2022a{]} and replay-based anti-forgetting {[}Robins, 1995{]} ---
are not novel. Our contribution is the empirical study and its honest
evaluation in a previously unexamined setting.

\hypertarget{related-work}{%
\subsection{2. Related Work}\label{related-work}}

\textbf{Domain adaptation for NMT.} Fine-tuning a general model on
in-domain data is the standard approach {[}Luong and Manning, 2015;
Freitag and Al-Onaizan, 2016{]}, surveyed by Chu and Wang {[}2018{]}.
Its central failure mode is catastrophic forgetting of the general
domain; mixing in general data during fine-tuning (``mixed
fine-tuning'') is a known mitigation {[}Chu et al., 2017{]}, a form of
experience replay {[}Robins, 1995{]}.

\textbf{Model soups / weight averaging.} Wortsman et al.~{[}2022a{]}
show that averaging the weights of multiple fine-tuned models improves
accuracy; WiSE-FT {[}Wortsman et al., 2022b{]} averages a fine-tuned
model with its zero-shot initialization to trade off in- and
out-of-distribution performance. We apply WiSE-FT-style base$\leftrightarrow$fine-tune
interpolation to MT domain adaptation.

\textbf{Indic MT.} IndicTrans2 {[}Gala et al., 2023{]} and BPCC are the
base model and a data source here. Prior WAT editions {[}e.g., Nakazawa
et al., 2024{]} host Indic MT shared tasks.

\textbf{Evaluation and its limits.} We report chrF2 {[}Popović, 2015{]}
via sacreBLEU {[}Post, 2018{]}, with paired bootstrap resampling for
significance {[}Koehn, 2004{]}; FLORES-200 {[}NLLB Team, 2022{]} is our
neutral general-domain test. Character-n-gram metrics are standard and
reproducible, but they measure overlap with references rather than
quality directly, and can diverge from human judgment {[}Callison-Burch
et al., 2006; Mathur et al., 2020; Freitag et al., 2021{]} --- a caveat
we take seriously and discuss in §6 and §8.

\hypertarget{data}{%
\subsection{3. Data}\label{data}}

\textbf{Base model.} \texttt{ai4bharat/indictrans2-en-indic-1B} (1.1B
parameters).

\textbf{Conversational data} (English$\to$Indic, each pair tagged with its
target language): - OPUS \textbf{OpenSubtitles} {[}Lison and Tiedemann,
2016{]}: en$\to$\{hi, ml, ta, te, ur\}, the only Indic languages with
substantial Indic subtitle data; capped at 40k/language to limit
dominance. - \textbf{BPCC-H-Daily} {[}Gala et al., 2023{]}:
human-authored everyday/assistant utterances; available for all 21
languages, \textasciitilde1.8k--11k pairs each. - \textbf{Tatoeba}:
everyday sentences (Hindi pool).

After filtering (length, length-ratio, exact dedup, Unicode-script
language identification), the conversational set is \textbf{294,119
pairs} across 21 languages (train 283,753 / dev 4,183 / test 6,183),
split per language. The Hindi-only experiments additionally apply LaBSE
cosine $\ge$ 0.70 to remove misaligned subtitle pairs.

\textbf{General anchor} (for experience replay): \textbf{BPCC-H-Wiki}
{[}Gala et al., 2023{]}, a general/encyclopedic register, sampled $\approx$1:1
per language to the conversational counts (254,469 pairs). The mixed
training set is therefore \textbf{538,222 pairs} (283,753 conversational
+ 254,469 general).

Per-language conversational volume is highly skewed: the five
OpenSubtitles languages (Hindi, Malayalam, Tamil, Telugu, Urdu) hold
30k--48k pairs each; the remaining sixteen rely on BPCC-H-Daily
($\approx$2k--11k). This skew is intrinsic to the available data and shapes our
results (§6).

\hypertarget{method}{%
\subsection{4. Method}\label{method}}

We compare three configurations, all fine-tuning IndicTrans2-1B
(IndicProcessor preprocessing with per-example target-language tags;
bf16; Adafactor; single NVIDIA A10G GPU):

\begin{enumerate}
\def\labelenumi{\arabic{enumi}.}
\tightlist
\item
  \textbf{Naive FT} --- fine-tune on conversational data only.
\item
  \textbf{Mixed FT (experience replay)} --- fine-tune on conversational
  + general anchor ($\approx$1:1).
\item
  \textbf{Model soup} --- weight-average the base and a fine-tuned
  model, $\theta$ = (1 $-$ $\alpha$)$\cdot$$\theta$(base) + $\alpha$$\cdot$$\theta$(FT), sweeping $\alpha$ and selecting on the
  conversational/general trade-off. Averaging is done by streaming the
  two checkpoints' tensors to keep memory bounded.
\end{enumerate}

\textbf{Training.} Hindi-only models use 146k conversational pairs (2--3
epochs); the multilingual models use the full 294k conversational / 538k
mixed sets (2 epochs / 1 epoch respectively, matched in total examples
seen).

\textbf{Evaluation.} chrF2 (sacreBLEU) on (a) a held-out conversational
test set and (b) FLORES-200 devtest (general), per language. We report
paired bootstrap significance (1000 resamples) for representative
languages.

\hypertarget{results}{%
\subsection{5. Results}\label{results}}

\hypertarget{single-language-case-study-hindi}{%
\subsubsection{5.1 Single-language case study
(Hindi)}\label{single-language-case-study-hindi}}

Table 1 isolates the mechanism on Hindi (conversational test =
subtitle-heavy held-out set; FLORES = 997-sentence devtest).

\textbf{Table 1: Hindi (chrF2).}

\begin{longtable}[]{@{}lll@{}}
\toprule\noalign{}
Model & conv & FLORES \\
\midrule\noalign{}
\endhead
\bottomrule\noalign{}
\endlastfoot
IndicTrans2 base & 51.6 & 61.7 \\
Naive conv-FT & \textbf{57.0} & 57.8 $\downarrow$ (forgetting) \\
Mixed FT & 56.0 & 60.6 \\
\textbf{Model soup ($\alpha$=0.6)} & 54.1 & \textbf{61.8} \\
\end{longtable}

Naive fine-tuning gains +5.4 chrF on conversational but loses $-$3.9 on
FLORES --- clear catastrophic forgetting. Mixed FT recovers most of the
general quality; the soup (mixed-FT averaged with base) is \textbf{$\ge$
base on both axes}. An $\alpha$-sweep (0.3--0.7) shows $\alpha$=0.6 is the knee: below
it general quality saturates at the base, above it FLORES drops below
base.

\hypertarget{all-21-languages-mixed-soup-vs.-base}{%
\subsubsection{5.2 All 21 languages (mixed-soup
vs.~base)}\label{all-21-languages-mixed-soup-vs.-base}}

Table 2 gives every language (n=200 sentences per test). The mixed-soup
model improves conversational chrF in \textbf{all 21 languages} (mean
\textbf{+6.2}) and changes FLORES by a mean of \textbf{$-$0.17} (all
within $\pm$0.7).

\textbf{Table 2: All 21 languages, base $\to$ mixed-soup (chrF2).}

\begin{longtable}[]{@{}lll@{}}
\toprule\noalign{}
Lang & conv base$\to$soup ($\Delta$) & FLORES base$\to$soup ($\Delta$) \\
\midrule\noalign{}
\endhead
\bottomrule\noalign{}
\endlastfoot
asm & 61.9$\to$69.4 (+7.6) & 44.2$\to$44.1 ($-$0.1) \\
ben & 68.7$\to$73.6 (+4.8) & 58.1$\to$58.1 ($-$0.0) \\
brx & 62.1$\to$68.0 (+5.9) & 45.9$\to$46.0 (+0.1) \\
doi & 63.4$\to$71.7 (+8.3) & 49.8$\to$49.6 ($-$0.2) \\
gom & 69.4$\to$75.6 (+6.2) & 49.1$\to$49.2 (+0.1) \\
guj & 53.7$\to$65.9 (+12.2) & 55.3$\to$54.7 ($-$0.6) \\
hin & 53.9$\to$56.8 (+2.9) & 60.0$\to$59.6 ($-$0.5) \\
kan & 52.4$\to$57.1 (+4.7) & 59.0$\to$58.7 ($-$0.3) \\
kas & 45.8$\to$51.3 (+5.5) & 39.1$\to$38.8 ($-$0.4) \\
mai & 57.8$\to$65.2 (+7.4) & 56.6$\to$56.6 ($-$0.0) \\
mal & 45.7$\to$48.4 (+2.7) & 61.1$\to$61.4 (+0.3) \\
mar & 70.2$\to$77.3 (+7.1) & 55.0$\to$54.5 ($-$0.5) \\
mni & 56.4$\to$62.8 (+6.4) & 48.2$\to$48.3 (+0.2) \\
npi & 68.8$\to$72.1 (+3.3) & 59.4$\to$59.0 ($-$0.5) \\
ory & 56.4$\to$65.7 (+9.2) & 52.0$\to$51.3 ($-$0.7) \\
pan & 68.0$\to$76.5 (+8.4) & 55.3$\to$54.8 ($-$0.5) \\
san & 54.5$\to$63.6 (+9.1) & 36.3$\to$36.5 (+0.2) \\
sat & 50.7$\to$56.2 (+5.6) & 30.7$\to$30.3 ($-$0.4) \\
tam & 52.7$\to$55.1 (+2.4) & 64.1$\to$64.4 (+0.3) \\
tel & 49.9$\to$60.0 (+10.1) & 62.4$\to$62.0 ($-$0.4) \\
urd & 48.1$\to$49.4 (+1.3) & 54.7$\to$54.7 (+0.1) \\
\textbf{mean} & \textbf{+6.2} & \textbf{$-$0.17} \\
\end{longtable}

\hypertarget{statistical-significance}{%
\subsubsection{5.3 Statistical
significance}\label{statistical-significance}}

We run paired bootstrap resampling {[}Koehn, 2004{]} (sacreBLEU, 1000
resamples, n=500) on the three languages with the hardest,
subtitle-based conversational tests (Table 3).

\textbf{Table 3: Paired bootstrap, base vs.~mixed-soup (chrF2; p = soup
vs.~base).}

\begin{longtable}[]{@{}lll@{}}
\toprule\noalign{}
Lang & conv (p) & FLORES (p) \\
\midrule\noalign{}
\endhead
\bottomrule\noalign{}
\endlastfoot
Hindi & 51.5 $\to$ 54.1 (\textbf{p=0.001}) & 61.9 $\to$ 61.9 (p=0.372, n.s.) \\
Tamil & 50.9 $\to$ 55.1 (\textbf{p=0.001}) & 64.7 $\to$ 64.6 (p=0.290, n.s.) \\
Malayalam & 46.9 $\to$ 49.5 (\textbf{p=0.004}) & 62.5 $\to$ 62.8
(\textbf{p=0.038}, $\uparrow$) \\
\end{longtable}

The conversational gains are significant (p $\le$ 0.004). FLORES shows
\textbf{no significant degradation} (not significant for Hindi/Tamil;
significantly \emph{improved} for Malayalam). The paired design is
important because conversational chrF has high variance (95\% CI $\approx$ $\pm$3
chrF at n=500), so unpaired CI overlap would be misleading; the paired
test isolates the per-sentence system difference.

\emph{(Absolute conversational chrF here differs from Table 2 ---
e.g.~Hindi 51.5 vs.~53.9 for the base --- because Table 3 uses n=500
sentences and Table 2 uses n=200; the high variance makes the exact mean
sample-dependent, which is precisely why we report the larger-n
\textbf{paired} test for the headline claims. The FLORES values, being
lower-variance, are stable across n.)}

\hypertarget{what-the-metric-does-and-does-not-establish}{%
\subsection{6. What the metric does and does not
establish}\label{what-the-metric-does-and-does-not-establish}}

Our claims are deliberately scoped to what chrF2 supports, because the
conversational setting has two specific ways the metric can mislead, and
we want to state them rather than paper over them.

\textbf{The general-domain claim is a tie, not a win.} On FLORES the
soup matches the base (Hindi 61.8 vs 61.7; all 21 languages within
$\pm$0.7). We therefore claim \emph{preservation} of general-domain quality,
not superiority on it.

\textbf{Part of the conversational gain is reference style-matching, not
necessarily adequacy.} chrF rewards character-n-gram overlap with the
reference. Our conversational references are subtitles, which are casual
in register; the adapted model is also more casual (informal pronouns,
spoken word order). A model that adopts the references' register
therefore gains chrF whether or not its \emph{meaning} is better.
Inspection of sentences where base and soup diverge shows exactly this
pattern --- the soup is consistently more colloquial --- so the
conversational chrF gain should be read as ``closer to the casual
reference style,'' which overlaps with but is not identical to ``a
better translation.'' This is the well-documented gap between n-gram
metrics and human judgment {[}Callison-Burch et al., 2006; Mathur et
al., 2020; Freitag et al., 2021{]}.

\textbf{What is solid.} Two things survive these caveats and are the
safe claims of the paper: (i) the recipe \emph{preserves} general-domain
chrF (FLORES tie, significance-tested), and (ii) it produces output
measurably closer to held-out human conversational references,
significantly so on the hardest subtitle tests. Whether (ii) is
perceived by humans as higher quality is a separate question that
automatic metrics cannot answer; see §8.

\hypertarget{analysis}{%
\subsection{7. Analysis}\label{analysis}}

\textbf{The forgetting is real, and the fix works.} Naive conversational
fine-tuning degrades FLORES ($-$3.9 for Hindi; a pure-conversational
multilingual soup degrades FLORES by $-$0.7 to $-$2.9 across languages).
Adding a general anchor before souping restores FLORES to base level
(mean $-$0.17), confirming experience replay as the operative mechanism.

\textbf{Gain magnitude is partly a test-set artifact.} The largest
conversational gains (Gujarati +12.2, Telugu +10.1, Odia +9.2) occur in
languages whose conversational \emph{test} is drawn solely from
BPCC-H-Daily --- short, formulaic sentences in the same style as
training. The most reliable gains are on languages whose test includes
harder, varied subtitle dialogue (Hindi +2.9, Tamil +2.4, Malayalam
+2.7) --- precisely the languages for which §5.3 confirms significance.
We therefore caution against ranking languages by raw $\Delta$conv across
heterogeneous test registers.

\textbf{A negative result on low-resource languages.} Languages with
very low absolute base quality --- Santali (FLORES 30.7), Sanskrit
(36.3), Kashmiri (39.1), Bodo (45.9) --- remain low after adaptation.
Their \emph{general} quality is bounded by the base model, which is in
turn bounded by the scarcity of any parallel data for these languages.
Conversational adaptation improves their conversational chrF but cannot
lift their general ceiling; doing so would require additional
\emph{general} parallel data, which does not exist publicly at
sufficient scale. This is a limitation of the data, not the method.

\hypertarget{conclusion}{%
\subsection{8. Conclusion}\label{conclusion}}

Using experience replay and model souping --- two established techniques
--- we adapt IndicTrans2 to conversational register across all 21 Indic
languages while preserving general-domain quality, with statistically
significant conversational chrF gains and no significant general-quality
loss. The recipe generalises across the family; its limits are set by
the base model's coverage of the lowest-resource languages. By automatic
metrics, the result is a single, drop-in model that \textbf{beats the
base on conversational chrF} (significantly; p $\le$ 0.004) while
\textbf{matching it on general-domain FLORES} --- a strictly
Pareto-non-worse trade on these two metrics. We are careful to keep this
claim scoped: it is a chrF result on a conversational test set, not an
unqualified quality claim. Confirming that the conversational chrF gain
is perceived by humans as higher quality --- as opposed to a register
shift the metric rewards (§6) --- is the natural next step (§9).

\hypertarget{limitations}{%
\subsection{9. Limitations}\label{limitations}}

\begin{itemize}
\tightlist
\item
  \textbf{Automatic metrics only.} All quality numbers are chrF2.
  Character-n-gram metrics measure overlap with references, not quality
  directly, and can overstate domain-adaptation gains when the in-domain
  references have a distinctive style: a model that adopts that style
  gains overlap independent of adequacy {[}Callison-Burch et al., 2006;
  Mathur et al., 2020; Freitag et al., 2021{]}. Our conversational
  references are subtitles, so part of the conversational chrF gain
  plausibly reflects register-matching (§6). Confirming a
  human-\emph{perceived} quality improvement requires a controlled,
  multi-annotator human evaluation; this is our priority next step and
  is \textbf{not yet conclusive}, so we make no human-quality claim and
  scope the contribution to the metric.
\item
  \textbf{Single training seed.} We report one run per configuration;
  training-variance across seeds is not yet quantified (significance
  §5.3 addresses test-set variance, not training variance).
\item
  \textbf{Significance on a subset.} Paired bootstrap is reported for 3
  of 21 languages (the hardest-test ones); extending to all 21 is
  straightforward future work.
\item
  \textbf{Register confound across languages.} Conversational test
  register differs by language (subtitles vs.~daily utterances),
  confounding cross-language gain comparison (§7).
\item
  \textbf{One direction.} English$\to$Indic only; Indic$\to$English and
  Indic$\leftrightarrow$Indic are untested and could exhibit different forgetting
  behaviour.
\end{itemize}

\hypertarget{reproducibility}{%
\subsection{Reproducibility}\label{reproducibility}}

All code is in the accompanying repository: data construction
(\texttt{build\_multiling\_conv.py}, \texttt{build\_ml\_mixed.py}),
training (\texttt{finetune\_it2\_ml.py}), souping
(\texttt{make\_soup.py}), evaluation (\texttt{ml\_eval.py}), and
significance (\texttt{sig\_test.py}). Released artifacts: the mixed-FT
model (\texttt{ft\_it2\_ml\_mixed}) and its soup
(\texttt{ft\_it2\_ml\_mixed\_soup}). Environment: Python 3.11,
\texttt{transformers==4.40.2}, \texttt{IndicTransToolkit}, sacreBLEU
2.6.

\hypertarget{references}{%
\subsection{References}\label{references}}

\begin{itemize}
\tightlist
\item
  Callison-Burch, Osborne, Koehn. \emph{Re-evaluating the Role of BLEU
  in MT Research.} EACL 2006.
\item
  Chu, Dabre, Kurohashi. \emph{An Empirical Comparison of Domain
  Adaptation Methods for NMT.} ACL 2017.
\item
  Chu and Wang. \emph{A Survey of Domain Adaptation for NMT.} COLING
  2018.
\item
  Freitag and Al-Onaizan. \emph{Fast Domain Adaptation for NMT.} arXiv
  2016.
\item
  Freitag, Foster, Grangier, et al.~\emph{Experts, Errors, and Context:
  A Large-Scale Study of Human Evaluation for MT.} TACL 2021.
\item
  French. \emph{Catastrophic Forgetting in Connectionist Networks.}
  Trends in Cognitive Sciences 1999.
\item
  Gala et al.~\emph{IndicTrans2: Towards High-Quality and Accessible MT
  for all 22 Scheduled Indian Languages.} TMLR 2023.
\item
  Goyal et al.~\emph{The FLORES-101 Evaluation Benchmark.} TACL 2022.
\item
  Koehn. \emph{Statistical Significance Tests for MT Evaluation.} EMNLP
  2004.
\item
  Lison and Tiedemann. \emph{OpenSubtitles2016.} LREC 2016.
\item
  Luong and Manning. \emph{Stanford NMT Systems for Spoken Language
  Domains.} IWSLT 2015.
\item
  Mathur, Baldwin, Cohn. \emph{Tangled up in BLEU: Reevaluating the
  Evaluation of Automatic MT Metrics.} ACL 2020.
\item
  McCloskey and Cohen. \emph{Catastrophic Interference in Connectionist
  Networks.} 1989.
\item
  Nakazawa et al.~\emph{Overview of the Workshop on Asian Translation.}
  WAT 2024.
\item
  NLLB Team (Costa-jussà et al.). \emph{No Language Left Behind.} arXiv
  2022. (FLORES-200)
\item
  Popović. \emph{chrF: Character n-gram F-score for Automatic MT
  Evaluation.} WMT 2015.
\item
  Post. \emph{A Call for Clarity in Reporting BLEU Scores} (sacreBLEU).
  WMT 2018.
\item
  Robins. \emph{Catastrophic Forgetting, Rehearsal and Pseudorehearsal.}
  Connection Science 1995.
\item
  Wortsman et al.~\emph{Model Soups.} ICML 2022a.
\item
  Wortsman et al.~\emph{Robust Fine-tuning of Zero-shot Models}
  (WiSE-FT). CVPR 2022b.
\end{itemize}

\end{document}
